\writeSectionFrame{\presentationTitle}{Model-View-Controller}
\begin{frame}{Steckbrief}
	\begin{itemize}
 		\item Deutscher Name: Model-View-Controller
 		\item Abkürzung: MVC
 		\item Englischer Name: Model-View-Controller
 		\item Weiterentwicklungen: Model2, MVP, MVVM
 		\item Art: Architekturmuster für Benutzeroberflächen
 		\item Im Grunde eine Familie von Mustern
 	\end{itemize}
\end{frame}

\begin{frame}{Grundlegende Idee}
	\begin{center}
 		\includegraphics[scale=0.25]{\BILDERPFAD/mvc1.png}
 	\end{center}
\end{frame}

\note{
	MVC unterteilt eine interaktive Anwendung in drei Bereiche: Verarbeitung, Ausgabe und Eingabe.	
}

\begin{frame}{Model}
	\begin{itemize}
 		\item Funktionaler Kern der Anwendung
 		\item Meldet abhängige Ansichten und Controller an
 		\item Benachrichtigt abhängige Komponenten bei Änderungen
 	\end{itemize}
\end{frame}

\begin{frame}{View}
	\begin{itemize}
 		\item Stellt die Informationen für den Anwender dar
 		\item Erzeugt und initialisiert ihre Controller
 		\item Implementiert die Aktualisierungsprozedur
 		\item Erfragt die Daten beim Modell
 	\end{itemize}
\end{frame}

\begin{frame}{Controller}
	\begin{itemize}
 		\item Erkennt Benutzereingaben als Ereignisse
 		\item Übersetzt Ereignisse in Dienstanforderungen an das Modell oder 
 		Darstellungsanforderungen an die Ansicht
 		\item Implementiert die Aktualisierungsprozedur wenn notwendig
 	\end{itemize}
\end{frame}

\begin{frame}{Aufbau}
	\begin{center}
 		\includegraphics[scale=0.25]{\BILDERPFAD/mvc2.png}
 	\end{center}
\end{frame}

\begin{frame}{Beispiel}
	\begin{itemize}
 		\item Einfaches Counter-MVC mit austauschbarem Modell und austauschbarer View
 		\item Models:
 			\begin{itemize}
 			  \item Model 1 erhöht einen Zähler um eins
 			  \item Model 2 verdoppelt den Zähler.
 			\end{itemize}
 		\item Views: Beide Views zeigen etwas andere Terminal-Steuerelemente und Texte an.
 		\item Controller: Kennt die Views und Models und koordiniert den Ablauf
 		\item Dialog ermöglicht es, das Programm mit View und Model zu konfigurieren
 	\end{itemize}
\end{frame}

\begin{frame}[fragile]{Model-Interface}
	\begin{exampleblock}{CounterModel.java}
 		\begin{lstlisting}
public interface CounterModel {
	int getCount();
	void increment();
}
 		\end{lstlisting}
 	\end{exampleblock}
\end{frame}

\begin{frame}[fragile]{Model-Implementierung}
	\begin{exampleblock}{IncrementCounterModel.java}
 		\begin{lstlisting}
public class IncrementCounterModel implements CounterModel {
    private int count = 0;
    private CounterChangeListener onChange;

    @Override
    public int getCount() {
        return count;
    }

    @Override
    public void increment() {
        count++;
        if (onChange != null) { 
        	onChange.onCountChanged(count);
        }
    }
}
 		\end{lstlisting}
 	\end{exampleblock}
\end{frame}

\begin{frame}[fragile]{View-Interface}
	\begin{exampleblock}{CounterView.java}
 		\begin{lstlisting}
public interface CounterView {
	void display(int count);
	void showMenu();
	boolean checkQuit(String input);
	boolean checkIncrement(String input);
	void printQuit();
	String readUserInput();
	void printGreeting();
}
 		\end{lstlisting}
 	\end{exampleblock}
\end{frame}

\begin{frame}[fragile]{View-Interface}
	\begin{exampleblock}{CounterView1.java}
 		\begin{lstlisting}
public class CounterView1 implements CounterView {
	@Override
	public void printGreeting() {
		System.out.println("Willkommen zu View1");
	}
	
	@Override
	public boolean checkIncrement(String input) {
		if (input.equals(INCREMENT)) {
			return true;
		} else {
			return false;
		}
	}
	
	@Override
	public String readUserInput() {
		String input = scanner.nextLine();
		return input;
	}
}
 		\end{lstlisting}
 	\end{exampleblock}
\end{frame}

\begin{frame}[fragile]{View-Interface}
	\begin{exampleblock}{CounterController.java}
 		\begin{lstlisting}
public class CounterController {
    private CounterModel model;
    private CounterView view;

    public CounterController(CounterModel model,
    		 CounterView view) {
        this.model = model;
        this.view = view;
        model.setChangeListener(
        	newCount -> view.display(newCount)); 
    }
    
    public void setModel(CounterModel model) {
		this.model = model;
	}
}
 		\end{lstlisting}
 	\end{exampleblock}
\end{frame}

\begin{frame}[fragile]{View-Interface}
	\begin{exampleblock}{CounterController.java}
 		\begin{lstlisting}
public class CounterController {
    public void run() {
        view.printGreeting();
        view.display(model.getCount());

        String input = "";
        while (view.checkQuit(input) == false) {
            view.showMenu();
            input = view.readUserInput();
            if (view.checkIncrement(input)) {
                model.increment();
            } 
        }
        
        view.printQuit();
    }
}
 		\end{lstlisting}
 	\end{exampleblock}
\end{frame}

%\begin{frame}{Ablauf einer Datenänderung}
%	\begin{center}
% 		\includegraphics[scale=0.25]{\BILDERPFAD/mvc6.png}
% 	\end{center}
%\end{frame}

\begin{frame}{Kombination mit Beobachtermuster}
	\begin{center}
 		\includegraphics[scale=0.3]{\BILDERPFAD/mvc7.png}
 	\end{center}
\end{frame}
\note{
	Das Beobachter-Muster ist der zentrale Mechanismus für die 
	Kommunikation zwischen Model und View in MVC. Es kann in verscheidenen
	Varianten implementiert werden. 
}

\begin{frame}{Model-View-Presenter oder Model-View-ViewModel}
	\begin{itemize}
	  \item Model hält eine Liste von Beobachtern
	  \item Es bietet Methoden an, um Beobachter hinzuzufügen und ggf. zu entfernen.
	  \item Wenn sich sein Zustand ändert, benachrichtigt es alle registrierten Beobachter
	  \item Jede View, die am Zustand des Models interessiert ist, implementiert die Schnittstelle des Beobachters
	\end{itemize}
\end{frame}

\begin{frame}[fragile]{View-Interface}
	\begin{exampleblock}{CounterChangeListener.java}
 		\begin{lstlisting}
@FunctionalInterface
public interface CounterChangeListener {
	void onCountChanged(int newCount);
}
 		\end{lstlisting}
 	\end{exampleblock}
\end{frame}

\begin{frame}[fragile]{View-Interface}
	\begin{exampleblock}{DoubleCounterModel.java}
 		\begin{lstlisting}
public class DoubleCounterModel implements CounterModel {
    private int count = 1;
    private CounterChangeListener onChange;

    // ...
	
	@Override
	public void setChangeListener(CounterChangeListener listener) {
		this.onChange = listener;
	}
}
 		\end{lstlisting}
 	\end{exampleblock}
\end{frame}

\begin{frame}[fragile]{Controller}
	\begin{exampleblock}{CounterController.java}
 		\begin{lstlisting}
public class CounterController {
    private CounterModel model;
    private CounterView view;

    public CounterController(CounterModel model, 
    		CounterView view) {
        model.setChangeListener(
        	newCount -> view.display(newCount)); 
    }
    
    public void setView(CounterView view) {
		this.view = view;
		model.setChangeListener(
			newCount -> view.display(newCount)); 
	}
    // ...
}
 		\end{lstlisting}
 	\end{exampleblock}
\end{frame}
\note{
	Der Controller stellt die Verbindung zwischen View und Model her. Er 
	registriert sich \glqq im Geiste{\grqq} als der Beobachter für das Model 
	und delegiert die Aktualisierung dann an die View. Die View ist der 
	Endempfänger der Aktualisierung.
}

\begin{frame}{Passive View oder Supervising Controller}
	\begin{itemize}
	  \item Der Controller kann isoliert getestet werden
	  \item View muss nichts über das Model wissen, außer welche Daten sie anzeigen soll
	  \item View hat keine Abhängigkeit zum Model.
	  \item In Varianten auch als 
	\end{itemize}
\end{frame}

\begin{frame}{Kombination mit Kompositummuster}
	\begin{center}
 		\includegraphics[scale=0.25]{\BILDERPFAD/mvc8.png}
 	\end{center}
\end{frame}

\begin{frame}{Kombination mit Strategiemuster}
	\begin{center}
 		\includegraphics[scale=0.25]{\BILDERPFAD/mvc9.png}
 	\end{center}
\end{frame}

\begin{frame}{Kombination mit anderen Mustern}
	\begin{center}
 		\includegraphics[scale=0.2]{\BILDERPFAD/mvc10.png}
 	\end{center}
\end{frame}

\begin{frame}{MVC als Musterfamilie}
	\begin{itemize}
	  \item MVC hat sich im Laufe der Zeit weiterentwickelt
	  \item Es gibt verschiedene Varianten für verschiedene Zwecke
	  \item Kernprinzip ist gleich: Trennung der Verantwortlichenkeiten zwischen Model, View und Controller 
	\end{itemize}
\end{frame}

\begin{frame}{Der Rest der Familie}
	\begin{itemize}
	  \item Klassisches Smalltalk-MVC
	  \item Document/View
	  \item Model-View-Presenter (MVP)
	  \item Model-View-ViewModel (MVVM)
	  \item Model2
	\end{itemize}
\end{frame}

\begin{frame}{Klassisches Smalltalk-MVC}
	\begin{itemize}
	  \item Entwickelt in den 70ern/80ern für Smalltalk
	  \item Relevant für historisches Verständnis
	\end{itemize}
	\pause
	\begin{block}{Model}
		\begin{itemize}
		  \item Enthält die Anwendungslogik und Daten 
		  \item Ist völlig unabhängig von View und Controller
		  \item Benachrichtigt Beobachter über Änderungen
		\end{itemize}
	\end{block}
\end{frame}

\begin{frame}{Klassisches Smalltalk-MVC}
	\begin{block}{View}
		\begin{itemize}
		  \item Zeigt den Zustand des Models an
		  \item Registriert sich selbst als Beobachter beim Model
		  \item Holt sich Daten direkt vom Model
		  \item Sendet Benutzereingaben als Events an den Controller. 
		\end{itemize}
	\end{block}
	\pause
	\begin{block}{Controller}
		\begin{itemize}
		  \item Empfängt Benutzereingaben von der View
		  \item Interpretiert diese Events und aktualisiert das Model oder die View
		  \item Reagiert auch auf Änderungen im Model, indem er die View anweist, sich zu aktualisieren 
		\end{itemize}
	\end{block}
\end{frame}

\begin{frame}{Klassisches Smalltalk-MVC}
	\begin{exampleblock}{Vorteile}
		\begin{itemize}
		  \item Klare Trennung der Verantwortlichkeiten
		  \item Hohe Wiederverwendbarkeit von Model und View.
		\end{itemize}
	\end{exampleblock}
	\pause
	\begin{alertblock}{Nachteile}
		\begin{itemize}
		  \item Starke Kopplung zwischen View und Controller (oft ein 1:1-Paar)
		  \item Die View muss das Model kennen, um sich zu aktualisieren
		  \item Der Controller muss oft wissen, wie die View die Events erzeugt
		  \item Für komplexe GUIs schwierig 
		\end{itemize}
	\end{alertblock}
\end{frame}

% -------------------

\begin{frame}{Document/View}
	Häufig in Frameworks wie Microsoft Foundation Classes (MFC) oder Java AWT/Swing
	\pause
	\begin{block}{Model / Document}
		Enthält die Daten
	\end{block}
	\begin{block}{View}
		\begin{itemize}
		  \item Eine oder mehrere Views zeigen Teile des Dokuments an
		  \item Sie sind meist direkt für die Verarbeitung von Benutzereingaben verantwortlich
		  \item Controller-Logik ist oft in der View
		\end{itemize}
	\end{block}	
\end{frame}

\begin{frame}{Document/View}
	\begin{exampleblock}{Vorteile}
		Einfacher für Anwendungen, die ein zentrales \glqq Dokument{\grqq} bearbeiten
	\end{exampleblock}
	\pause
	\begin{alertblock}{Nachteile}
		\begin{itemize}
		  \item Die Trennung von View und Controller ist oft schwach oder existiert nicht
		  \item Wiederverwendbarkeit der View ist erschwert
		  \item Der \glqq Controller{\grqq} ist oft nur Event-Listener in der View.
		\end{itemize}
	\end{alertblock}
\end{frame}

% ----------------
\begin{frame}{Model-View-Presenter (MVP)}
	\begin{itemize}
	  \item Entstanden wegen Testens von GUIs
	  \item Trennung sollte verbessert werden
	  \item Besonders populär in Windows Forms, ASP.NET Web Forms, Android-Entwicklung.
	\end{itemize}
	\pause
	\begin{block}{Model}
		Wie im klassischen MVC
	\end{block}
\end{frame}

\begin{frame}{Model-View-Presenter (MVP)}
	\begin{block}{View}
		\begin{itemize}
		  \item Eine passive View (Dumb View)
		  \item Zeigt Daten an, die ihr der Presenter gibt
		  \item Sammelt Benutzereingaben und leitet sie sofort an den Presenter weiter
		  \item Hat keine Logik, keine Kenntnis vom Model. 
		  \item Definiert eine Schnittstelle (Interface), über die der Presenter mit ihr kommuniziert. 
		\end{itemize}
	\end{block}
\end{frame}

\begin{frame}{Model-View-Presenter (MVP)}
	\begin{block}{Presenter}
		\begin{itemize}
		  \item Agiert als der Controller und die Logik für die View
		  \item Empfängt Eingaben von der View
		  \item Aktualisiert das Model und aktualisiert dann die View direkt über ihre Schnittstelle
		  \item Hat volle Kontrolle über die View 
		\end{itemize}
	\end{block}
\end{frame}

\begin{frame}{Model-View-Presenter (MVP)}
	\begin{exampleblock}{Vorteile}
		\begin{itemize}
		  \item Extrem hohe Testbarkeit (Presenter ist testbar ohne GUI)
		  \item View ist komplett entkoppelt vom Model
		  \item View ist sehr gut wiederverwendbar.
		\end{itemize}
	\end{exampleblock}
	\pause
	\begin{alertblock}{Nachteile}
		 Presenter kann sehr dick werden (Fat Presenter), da er viel UI-Logik 
		 enthält.
	\end{alertblock}
\end{frame}

% ----------------
\begin{frame}{Model-View-ViewModel (MVVM)}
	\begin{itemize}
	  \item Entwickelt von Microsoft für WPF und Silverlight 
	  \item Später übernommen von anderen Frameworks
	\end{itemize}
	\pause
	\begin{block}{Model}
		Wie im klassischen MVC
	\end{block}
\end{frame}

\begin{frame}{Model-View-Presenter (MVP)}
	\begin{block}{View}
		\begin{itemize}
		  \item Bindet (Data Binding) direkt an Eigenschaften und Befehle des ViewModels
		  \item Hat wenig bis keine Logik
		  \item Muss das ViewModel kennen, aber das ViewModel kennt die View nicht direkt 
		\end{itemize}
	\end{block}
\end{frame}

\begin{frame}{Model-View-Presenter (MVP)}
	\begin{block}{ViewModel}
		\begin{itemize}
		  \item Abstraktion der View
		  \item Enthält die Daten und Befehle, die die View benötigt, in einem für die View geeigneten Format
		  \item Konvertiert Model-Daten für die View und View-Eingaben für das Model
		  \item Bindet an das Model und aktualisiert sich bei dessen Änderungen
		  \item Implementiert oft das Observer-Pattern, um Änderungen an die View zu melden
		  \item Hat keine direkte Referenz zur View, Kommunikation erfolgt über Data Binding. 
		\end{itemize}
	\end{block}
\end{frame}

\begin{frame}{Model-View-Presenter (MVP)}
	\begin{exampleblock}{Vorteile}
		\begin{itemize}
		  \item Sehr hohe Testbarkeit (ViewModel ist testbar ohne GUI)
		  \item View und ViewModel sind stark entkoppelt
		  \item Ermöglicht die Arbeit von UI-Designern und Entwicklern parallel.
		\end{itemize} 
	\end{exampleblock}
	\pause
	\begin{alertblock}{Nachteile}
		\begin{itemize}
		  \item Data Binding kann komplex werden
		  \item Performance-Probleme bei sehr vielen Bindings
		  \item Hoher Initialaufwand für das Setup.
		\end{itemize}
	\end{alertblock}
\end{frame}

%-----------------

\begin{frame}{Model2}
	\begin{itemize}
	  \item Ursprung Ende der 1990er, Anfang der 2000er Jahre
	  \item Geprägt durch Java EE (insbesondere JavaServer Pages - JSP und Servlets),
	  \item Reaktion auf die Probleme von Model 1-Architekturen, bei denen JSPs direkt die Geschäftslogik und Datenbankzugriffe enthielten
	  \item Im Wesentlichen eine serverseitige Adaption des MVC-Musters für das Request-Response-Paradigma des Webs
	  \item eine spezielle Form von MVC, die die Besonderheiten des Web (stateless HTTP, Request/Response-Zyklus) berücksichtigt
	  \item Es betont eine passive View (ähnlich MVP) und einen aktiven Controller, der die gesamte Flusskontrolle übernimmt.
	\end{itemize}
\end{frame}

\begin{frame}{Model-View-Presenter (MVP)}
	\begin{block}{Model}
		\begin{itemize}
		  \item Repräsentiert die Geschäftslogik, Daten und den Anwendungszustand
		  \item Ist völlig unabhängig von der Web-spezifischen Präsentation.
		\end{itemize}
	\end{block}
	\pause
	\begin{block}{View}
		\begin{itemize}
		  \item Verantwortlich für die Generierung der Ausgabeseite (HTML, XML, JSON etc.), die an den Browser gesendet wird
		  \item Ist passiv und hat keine Geschäftslogik.
		  \item Empfängt die Daten, die es anzeigen soll, vom Controller 
		\end{itemize}
	\end{block}
\end{frame}

\begin{frame}{Model-View-Presenter (MVP)}
	\begin{block}{Controller}
		\begin{itemize}
		  \item Herzstück des Model 2
		  \item Der einzige Endpunkt, der HTTP-Anfragen vom Client empfängt
		  \item Keine Präsentationslogik und keine Datenzugriffslogik
		  \item Aufgaben:
		  	\begin{itemize}
			  \item Analyse der eingehenden HTTP-Anfrage
			  \item Delegation an das Model zur Ausführung der Geschäftslogik.
			  \item Empfängt die Ergebnisse vom Model.
			  \item Wählt die entsprechende View zur Darstellung aus
			  \item Bereitet die Daten für die View vo
			  \item Leitet die Anfrage an die ausgewählte View weiter 
			\end{itemize}
		\end{itemize}
	\end{block}
\end{frame}

\begin{frame}{Model-View-Presenter (MVP)}
	\begin{exampleblock}{Vorteile}
		\begin{itemize}
		  \item Sehr hohe Testbarkeit (ViewModel ist testbar ohne GUI)
		  \item View und ViewModel sind stark entkoppelt
		  \item Ermöglicht die Arbeit von UI-Designern und Entwicklern parallel.
		\end{itemize} 
	\end{exampleblock}
	\pause
	\begin{alertblock}{Nachteile}
		\begin{itemize}
		  \item Data Binding kann komplex werden
		  \item Performance-Probleme bei sehr vielen Bindings
		  \item Hoher Initialaufwand für das Setup.
		\end{itemize}
	\end{alertblock}
\end{frame}
